% Section 2 -- Setting and background. Flowing moderate sentences, no clause-stacked lists, no
% forward-pointing signposts, no casual/defensive phrasing. Verified citations only.
\section{Setting and background}
\label{sec:setting}

We forecast in-game win probability in American football. The state at a given play comprises the score
margin, the quarter and time remaining, down and distance, field position, the team in possession, and
the public pregame point spread. The forecaster returns a single probability that the team in possession
wins. We draw states from regular-season National Football League games and partition them by season,
training on 2015 through 2022, selecting on 2023, and testing on 2024, so that no game appears in more
than one split.
The only market signal in the prompt is the pregame point spread, which is fixed before kickoff and
public; the live win probability the market quotes during the game is withheld from both the model and
the reward, and enters only at evaluation.

The output is a probability, and the label is a single realized outcome of a stochastic event. This
separates the task from the calibration problems usually studied in reinforcement learning, where the
model answers a question that has a correct answer and the quantity of interest is its confidence that
the answer is right \citep{rlcr, dcpo}. That uncertainty is epistemic, and a stronger model can in
principle reduce it. The uncertainty here is aleatoric: the outcome is random given the state, and no
forecaster can remove it. The target of calibration is therefore the conditional win rate
$\eta(x)=\Pr(\text{win}\mid x)$, not the correctness of an answer.

We measure a forecaster by how closely its probabilities match observed frequencies. The Brier score,
the mean squared error $(p-y)^2$ between a forecast $p$ and an outcome $y\in\{0,1\}$ \citep{brier1950},
is strictly proper: among all functions of the state, its expectation is minimized uniquely by $\eta(x)$
\citep{gneiting2007}. It decomposes into reliability, resolution, and uncertainty \citep{murphy1973}.
Reliability is calibration, the agreement between stated probabilities and realized frequencies;
resolution is sharpness, the spread of forecasts across states of differing outcome rate; uncertainty is
fixed by the base rate. Because two forecasters can be equally calibrated yet differ in Brier score
through resolution alone, we report calibration on its own, through the expected and maximum calibration
error and reliability diagrams \citep{guo2017, brocker2007, dimitriadis2021}. We compare forecasters with
a paired bootstrap over plays \citep{efron1979}. Calibration error also tests whether training found
$\eta$ rather than memorizing the training outcomes, since a model that memorized them would be
miscalibrated on held-out games.

We take the betting market as the ceiling: it quotes a live win probability for each state and predicts
outcomes more accurately than statistical models built from game features \citep{boulier2003,
levitt2004}. We convert its odds to a probability \citep{strumbelj2014}. The gap between a forecaster and
the market then measures the live in-game information that the public state lacks.
